\chapter{群论}

% ===================== 第一节 =====================
\section{关系、等价关系与分类}

\begin{definition}[二元关系]
设 $A$ 为集合，$D=\{\text{对},\text{错}\}$，那么 $A\times A$ 到 $D$ 的每个映射 $R$ 就叫做 $A$ 的一个关系（也称为二元关系）。

若 $R:(a,b)\to\text{对}$，就称 $a$ 与 $b$ 符合关系，记为 $aRb$；

若 $R:(a,b)\to\text{错}$，就称 $a$ 与 $b$ 不符合关系，记为 $a\overline{R}b$。

由上述定义知，$A$ 中任一对元 $a,b$，都可以判定是否符合这个关系。
\end{definition}

\begin{note}[关系的定义 二元关系]
二元关系 每一个映射都叫做A的一个关系 D中对错是固定的集合 D不一定要 \\
R可以代表所有比较符号甚至定义的关系；R也可以是$A\times A$的一个子集；从定义上来将$aRb$的意思是$(a,b)$属于集合$R$
\end{note}

\begin{example}
集合 $A=\{1,2,3\}$ 上的"$\le$"关系：

$R=\{(a,b)\in A\times A\mid a\le b\}=\{(1,1),(1,2),(1,3),(2,2),(2,3),(3,3)\}$

$R$ 明确地被定义为 $A\times A$ 的一个子集。我们通常写 $aRb$ 来表示"$a$ 与 $b$ 满足关系 $R$"。
\end{example}

\begin{definition}[等价关系]
设 $\sim$ 是集合 $A$ 上的二元关系，如果 $\sim$ 具有以下三种性质：
\begin{enumerate}
  \item 反射律（反身性）：$\forall a\in A,\,a\sim a$；
  \item 对称律（对称性）：$\forall a,b\in A$，当 $a\sim b$ 时必有 $b\sim a$；
  \item 推移律（传递性）：$\forall a,b,c\in A$，当 $a\sim b$ 且 $b\sim c$ 时必有 $a\sim c$，
\end{enumerate}
则称 $\sim$ 叫做 $A$ 上的等价关系。并且当 $a\sim b$ 时，习惯称 $a$ 与 $b$ 等价。
\end{definition}

\begin{note}[等价关系]
$aRa$ / $aRb$则$bRa$ / $aRc$，$bRa$则$bRc$ \\
显然这里是R代表更广义的=除了等于的意思还表示$(a,c)$属于集合$R$ \\
比如大于号就不是等价关系
\end{note}

\begin{definition}[分类]
若把一个集合 $A$ 分成若干个叫做类的子集，使得 $A$ 的每一个元素属于而且只属于一个类，那么这些类的全体叫做集合 $A$ 的一个分类。

假定我们有一个集合的一个分类，那么，一个类里的任何一个元叫做这个类的一个\textbf{代表}。刚好由每一类的一个代表作成的集合叫做一个\textbf{全体代表团}。
\end{definition}

\begin{note}[分类]
1.$A=\bigcup A_i$ \quad 2.$A_i$非空 \quad 3.任意$A_i$和$A_j$交集是空 \\
分类是所有类的全体 / 类里面一个元素就是代表 / 一个子集叫一个类
\end{note}

\begin{theorem}[分类决定等价关系]
集合 $A$ 的每个分类都决定了 $A$ 的一个等价关系。

\textbf{证明：} 设 $\Omega=\{A_i,i\in I\}$ 是 $A$ 的一个一分类，用 $\Omega$ 我们可以规定 $A$ 上的一个二元关系 $\sim$：$a\sim b\Leftrightarrow a$ 和 $b$ 同在 $A_i$ 中。显然 $\sim$ 是 $A$ 的一个等价关系，须验证是等价关系：

（1）反射性：$\forall a\in A$，则 $a$ 在某个 $A_i$ 中，故 $a\sim a$；

（2）对称性：$\forall a,b\in A$，若 $a\sim b$，则 $a$ 和 $b$ 同在某个 $A_i$ 中，当然 $b\sim a$；

（3）传递性：$\forall a,b,c\in A$，若 $a\sim b$，$b\sim c$，放在 $A_i$ 中，则 $a,b,c$ 都同在 $A_i$ 中，所以 $a\sim c$。
\end{theorem}

\begin{note}[每一个分类都决定了A的一个等价关系]
产生一个分类就能产生一个等价关系 \\
因为一个分类里面都是有同一种性质的元素 \\
把这种性质作为等价的判定自然能够构造出一种等价关系
\end{note}

\begin{theorem}[等价关系决定分类]
集合 $A$ 上的任一个等价关系都可确定 $A$ 的一个分类。

\textbf{证明：} $\forall a\in A$，令 $[a]=\{x\in A\mid x\sim a\}$，如此确定的这些子集具有：

（1）$[a]\ne\varnothing$：由 $a\sim a\Rightarrow a\in[a]$；

（2）$[a]\cap[b]=\varnothing$，当 $a$ 与 $b$ 不等价时；若 $x\in[a]\cap[b]\Rightarrow x\sim a$，$x\sim b$，由 $\sim$ 的对称性和传递性知 $a\sim b$，推出矛盾，所以 $[a]\cap[b]=\varnothing$；

（3）$A=\bigcup_{a\in A}[a]$：$\because\forall a\in A\Rightarrow a\in[a]\subseteq\bigcup_{a\in A}[a]$。

$\therefore\Omega=\{[a]\mid\forall a\in A\}$ 是 $A$ 的一个分类。
\end{theorem}

\begin{note}[集合A的任一个等价关系都可确定A的一个分类]
条件给的是一个等价关系 要证明的是确定了一个分类 \\
2不同类的子集相交时 如果a集合和b集合不等价那么a集合和b集合的交集是空集
\end{note}

% ===================== 第二节 =====================
\section{映射}

\begin{definition}[映射]
映射定义要注意以下几点：
\begin{enumerate}
  \item 集合 $A_1,A_2,\cdots,A_n,D$ 可以相同；
  \item $A_1,A_2,\cdots,A_n$ 的次序不能掉换；
  \item 映射一定要替每一个元规定一个象；
  \item 一个元只能有唯一的象；
  \item 所有的象都必须是 $D$ 的元。
\end{enumerate}
\end{definition}

\begin{definition}[单射、满射与双射]
（3）双射（一一对应）：若 $f$ 既是单射又是满射，则 $f$ 是双射。

定义：一个 $A$ 到 $A$ 的映射叫做 $A$ 的一个变换。一个 $A$ 到 $A$ 的满射、单射或 $A$ 与 $A$ 间的一一映射叫做 $A$ 的一个满射变换、单射变换或一一变换。
\end{definition}

\begin{note}[单射满射]
单射就是原像不同 像不同 \\
映射 原像相等 像相等；像相等原像不一定相等 \\
单射 原像不等 像一定不等；像相等原像一定相等 \\
A到A的映射叫做A的一个变换
\end{note}

\begin{note}[逆映射]
如果$X\times Y$的一个子集是$X$到$Y$的一一映射$\varphi$，那么$\varphi$中的每个$x\in X$都作为第一个成员出现在一个有序对中，每个$y\in Y$都作为第二个成员出现在一个有序对中。因此，如果我们交换$\varphi$中所有有序对$(x,y)$的第一个成员和第二个成员，得到一个有序对$(y,x)$的集合，我们得到一个$Y\times X$的子集，它给出一个$Y$到$X$的一一映射。这个函数被称为$\varphi$的逆函数，用$\varphi^{-1}$表示。总结一下，如果$\varphi$将$X$一对一地映射到$Y$，且$\varphi(x)=y$，那么$\varphi^{-1}$将$Y$一对一地映射到$X$，且$\varphi^{-1}(y)=x$。
\end{note}

% ===================== 第三节 =====================
\section{同余关系与剩余类}

\begin{definition}[模 $n$ 的同余关系与剩余类]
任取 $0<n\in\mathbb{Z}$，可以在 $\mathbb{Z}$ 中确定一种等价关系
\[
R:\forall a,b\in\mathbb{Z},\;aRb\Leftrightarrow n\mid a-b
\]
则称 $R$ 为模 $n$ 的同余关系，并将 $aRb$ 记为 $a\equiv b\pmod{n}$（$a$ 同余 $b$ 模 $n$）。

称由同余关系确定的分类中的类为模 $n$ 的剩余类。
\end{definition}

\begin{note}[同余关系 模n的剩余类]
$a-b$能整除$n$说明$a/n$和$b/n$的余数相同
\end{note}

\begin{note}[剩余类加群]
剩余类加群有$A$必有$n-A$（有一个剩余类就有其对应的"负"类）。
\end{note}

% ===================== 第四节 =====================
\section{代数运算与单位元}

\begin{definition}[代数运算]
一个 $A\times B$ 到 $D$ 的映射叫做一个 $A\times B$ 到 $D$ 的代数运算。
\[
\circ:(a,b)\to d=\circ(a,b)
\]
$\circ(a,b)$ 可写成 $a\circ b$，$\therefore\;\circ:(a,b)\to d=a\circ b$。
\end{definition}

\begin{note}[代数运算定义]
两个集合的积到D的映射叫代数运算
\end{note}

\begin{example}
$A=\{\text{所有整数}\}$，$B=\{\text{所有不等于零的整数}\}$，$D=\{\text{所有有理数}\}$

$\circ:(a,b)\to\dfrac{a}{b}=a\circ b$

是一个 $A\times B$ 到 $D$ 的代数运算，也就是一个普通的除法。
\end{example}

\begin{definition}[封闭性]
若集合 $S$ 上的二元运算 $*$ 满足：对任意 $a,b\in S$，都有 $a*b\in S$，则称 $S$ 对运算 $*$ 是封闭的。
\end{definition}

\begin{note}[封闭 二元运算]
乘法与乘积 运算的所有元素都是A里面的元素 \\
实数集$+$、$-$、$\times$是二元运算 \\
二元运算也称为A上的代数运算
\end{note}

\begin{definition}[单位元]
设 $S$ 是一个具有二元运算 $*$ 的集合。如果对集合 $S$ 中的所有 $a$，$a*e=e*a=a$，那么 $e$ 被称为是二元运算 $*$ 的单位元（幺元）。

\textbf{定理（单位元唯一性）：} 具有二元运算 $*$ 的集合最多有一个单位元。
\end{definition}

\begin{note}[单位元唯一性证明]
单位元唯一性证明方法：把单位元拆解，然后用结合律引出新的东西，唯一性证明。

若 $e_1,e_2$ 都是单位元，则 $e_1=e_1*e_2=e_2$。

数学中唯一性证明通常设存在两个再证明这两个相等
\end{note}

\begin{definition}[逆元]
如果 $*$ 是 $S$ 上的一个运算，且有单位元 $e$，那么对任意 $x\in S$，$x$ 的逆元是 $x'$，使得 $x*x'=x'*x=e$。
\end{definition}

% ===================== 第五节 =====================
\section{幺半群}

\begin{definition}[幺半群（定义 1.1）]
我们说 $(S,*)$ 是一个幺半群，当该二元运算满足结合律，且具有单位元。此即，
\[
\forall x,y,z\in S,\;x*(y*z)=(x*y)*z
\]
\[
\exists e\in S,\forall x\in S,\;e*x=x*e=x
\]
\end{definition}

\begin{definition}[交换幺半群（定义 1.2）]
我们说 $(S,*)$ 是一个交换幺半群，当其是一个幺半群，且运算满足交换律，即
\[
\forall x,y\in S,\;x*y=y*x
\]
\end{definition}

\begin{definition}[幂（定义 1.3）]
令 $x_1,\cdots,x_n\in S$，我们递归地定义
\[
x_1\cdot x_2\cdots x_n=(x_1\cdot x_2\cdots x_{n-1})\cdot x_n
\]
令 $x\in S$，$n\in\mathbb{N}$。若 $n>0$，我们定义 $x^n=x\cdots x$，而 $x^0=e$。
\end{definition}

\begin{note}[定义1 3群中的幂]
牢记幂只是把多个乘法换了一种写法表示了，乘法只是一种运算
\end{note}

\begin{note}[幺半群的定义 定义1 2 定义1 3]
幺半群只要求有单位元，还有满足结合律 \\
幺半群的单位元是唯一的
\end{note}

\begin{proposition}[广义结合律（命题 1.2）]
令 $x_1,\cdots,x_n,y_1,\cdots,y_m\in S$，则
\[
x_1\cdot x_2\cdots x_n\cdot y_1\cdot y_2\cdots y_m=(x_1\cdot x_2\cdots x_n)\cdot(y_1\cdot y_2\cdots y_m)
\]
\end{proposition}

\begin{note}[广义结合律与数学归纳法]
数学归纳法先对1一个证明成立然后假设k成立证明k成立的时候可以借鉴证明特殊的做法再证明k+1成立
\end{note}

\begin{proposition}[命题 1.3]
令 $x\in S$，$m,n\in\mathbb{N}$，则 $x^{m+n}=x^m\cdot x^n$。
\end{proposition}

\begin{definition}[子幺半群（定义 1.4）]
令 $(S,\cdot)$ 是一个幺半群，若 $T\subset S$，我们说 $(T,\cdot)$ 是 $(S,\cdot)$ 的一个子幺半群，若 $e\in T$，且 $T$ 在乘法下封闭，即
\[
e\in T,\qquad\forall x,y\in T,\;x*y\in T
\]
\end{definition}

\begin{definition}[幺半群同态（定义 1.5）]
假设 $(S,\cdot)$，$(T,*)$ 是两个幺半群，且 $f:S\to T$ 是一个映射，我们称 $f$ 是一个幺半群同态，当 $f$ 保持了乘法运算，且把单位元映到了单位元。
\[
\forall x,y\in S,\;f(x\cdot y)=f(x)*f(y),\qquad f(e)=e'
\]
其中，$e$ 和 $e'$ 分别是 $(S,\cdot)$ 和 $(T,*)$ 的单位元。
\end{definition}

\begin{note}[定义1 5幺半群的同态]
这里是把自然数中的x映射到了x的n次方这个元素是在另一个幺半群里面的
\end{note}

\begin{definition}[由子集生成的子幺半群（定义 1.6）]
假设 $(S,\cdot)$ 是一个幺半群，而 $A\subset S$ 是一个子集。我们定义由 $A$ 生成的子幺半群，记作 $\langle A\rangle$，是指 $S$ 中所有包含了 $A$ 的子幺半群的交集。
\[
\langle A\rangle=\bigcap\{T\subset S:T\supset A,\,T\text{ 是子幺半群}\}
\]
\end{definition}

\begin{note}[定义1 6由子集生成的子幺半群]
是包含A的所有子幺半群，如果A本身就是幺半群它不需要添加项就可以生成子幺半群。但是如果A本身元素构成不了幺半群你们就要引入新的元素显然新的元素+A是包含A的，所有A生成的子幺半群里面都有全部的A。所有的交集也就A+最少构成子幺半群需要添加的元素，这个最少代价形成的幺半群就是A生成的子幺半群
\end{note}

\begin{definition}[幺半群同构（定义 1.7）]
假设 $(S,\cdot)$，$(T,*)$ 是两个幺半群，且 $f:S\to T$ 是一个映射，我们称 $f$ 是一个幺半群同构，当 $f$ 是双射，且是一个同态。
\[
f\text{ 是双射},\qquad\forall x,y\in S,\;f(x\cdot y)=f(x)*f(y),\qquad f(e)=e'
\]
\end{definition}

\begin{note}[定义1 7幺半群的同构]
同构概念可以类比全等 \\
双射，原函数同态
\end{note}

\begin{proposition}[命题 1.6]
若 $f:(S,\cdot)\to(T,*)$ 是一个幺半群同构，则 $f^{-1}:T\to S$ 是一个幺半群同构。
\end{proposition}

\begin{note}[命题1 5证明A生成的子幺半群是包含A的最小子幺半群]
由一个子集A生成的子结构就是包含了A的所有子结构的交集 \\
所有幺半群都有e自然在交集了吗，证明x，y有$x*y$就有因为每一个群x，y在交集里面所以每个半幺群里面都有
\end{note}

\begin{definition}[逆元（定义 1.8）]
令 $(S,\cdot)$ 是一个幺半群，$x\in S$。我们称 $x$ 是可逆的，当且仅当
\[
\exists y\in S,\;x\cdot y=y\cdot x=e
\]
其中 $y$ 被称为 $x$ 的逆元，记作 $x^{-1}$。
\end{definition}

\begin{lemma}[引理 1.1：幺半群的逆元构成群]
令 $(S,\cdot)$ 是一个幺半群，则 $G$ 是由所有可逆元构成的子集，则 $(G,\cdot)$ 是个群。
\end{lemma}

\begin{proposition}[命题 1.7：逆元唯一]
令 $(S,\cdot)$ 是一个幺半群。假设 $x\in S$ 是可逆的，则其逆元唯一。也就是说，如果 $y,y'\in S$ 都是它的逆元，则 $y=y'$。

\textbf{证明：} 假设 $y,y'$ 都是 $x$ 的逆元。则 $y\cdot x=e$，$x\cdot y'=e$。利用一个巧妙的连续"代数变形"来证明 $y=y'$。
\[
y=y\cdot e=y\cdot x\cdot y'=e\cdot y'=y'
\]
其中我们利用了 $e$ 是单位元，以及（广义）结合律。
\end{proposition}

% ===================== 第六节 =====================
\section{群}

\begin{definition}[群（定义 1.9）]
令 $(G,\cdot)$ 是一个幺半群，则我们说它是一个群，若 $G$ 中所有元素都是可逆的，换言之，若是 $G$ 上的一个二元运算，则我们称 $(G,\cdot)$ 是个群，若
\[
\forall x,y,z\in G,\;x\cdot(y\cdot z)=(x\cdot y)\cdot z
\]
\[
\exists e\in G,\forall x\in G,\;e\cdot x=x\cdot e=x
\]
\[
\forall x\in G,\;\exists y\in G,\;x\cdot y=y\cdot x=e
\]

注意：对于群，我们常用字母 $G$，以及接着的 $H,K$ 等，因为群的英文是 group。另外，我们把 $x$ 的逆记作 $x^{-1}$。命题 1.8：$(x^{-1})^{-1}=x$。命题 1.9：$(x\cdot y)^{-1}=y^{-1}\cdot x^{-1}$。
\end{definition}

\begin{note}[定义1 9幺半群加逆元就是群]
这里有着为什么加逆符号这个需要改变x，y的位置
\end{note}

\begin{note}
群的另一种等价定义：设 $G$ 是一个具有叫做乘法代数运算的非空集合，并且满足：
\begin{enumerate}
  \item[I.] $G$ 对于乘法来说是闭的；
  \item[II.] 结合律：$\forall a,b,c\in G$，有 $(a\circ b)\circ c=a\circ(b\circ c)$；
  \item[III.] $G$ 中有左单位元 $e$：$\forall a\in G$，$e\circ a=a$；
  \item[IV.] 对 $G$ 中每一个元素 $a$，有左逆元 $a^{-1}\in G$，使得 $a^{-1}\circ a=e$。
\end{enumerate}
则称 $G$ 关于代数运算 $\circ$ 作成一个群。
\end{note}

\begin{note}[群的定义]
乘法封闭 满足结合律 叫做乘法不一定是乘法是可以自定义的 \\
有左单位元有左逆元 \\
先证明是代数运算
\end{note}

\begin{note}[群的意义]
群论的诞生源于对具体数学问题的解决需求和对称性认识的深化：
\begin{enumerate}
  \item \textbf{多项式方程与伽罗瓦理论：} 伽罗瓦通过研究根的置换群，判断多项式是否能用有限次根式求解，奠定了群论基础。
  \item \textbf{几何中的对称性：} 旋转、反射、平移等几何变换具有规律性，群论统一描述这些对称操作，是几何学的强大工具。
  \item \textbf{数学结构的统一与抽象化：} 群的公理体系（封闭性、结合律、单位元、逆元）简单而有力，推动了环、域、向量空间等更复杂结构的建立，使数学理论趋向统一和抽象。
  \item \textbf{物理学的应用：} 物理定律在某些变换下保持不变（对称性与守恒定律），群论被用于量子力学、粒子物理等领域构建理论模型。
\end{enumerate}
群论是研究"对称性"与"结构"的核心工具，既解决具体数学问题，也统一描述各种抽象代数结构。
\end{note}

\begin{note}[群的性质]
$\forall a\in G,\exists a^{-1}\in G,\;a^{-1}a=e\Rightarrow\exists a'\in G,\;a'a^{-1}=e$

定理1：左逆元也是右逆元

$a$ 逆是 $a$ 的左逆元，$a'$ 是 $a$ 逆的左逆元，再证明 $a'$ 和 $a$ 一样

$aa^{-1}=eaa^{-1}=(a'a^{-1})aa^{-1}=a'(a^{-1}a)a^{-1}=a'ea^{-1}=a'a^{-1}=e$

左单位元也是右单位元；群的乘法适合消去律。

定理1左逆元也是右逆元 / $a$逆是$a$的左逆元 $a'$是$a$逆的左逆元再证明$a'$和$a$一样 \\
左单位元也是右单位元 / 群的乘法适合消去律
\end{note}

\begin{note}[群的等价判别]
设 $G$ 是一个具有叫做乘法代数运算的非空集合，并且满足：
\begin{enumerate}
  \item[I.] $G$ 对于乘法来说是闭的；
  \item[II.] 结合律：$\forall a,b,c\in G$，有 $(a\circ b)\circ c=a\circ(b\circ c)$；
  \item[III.] $G$ 中任意两个元 $a,b$，方程 $ax=b$，$ya=b$ 都在 $G$ 中有解。
\end{enumerate}
则称 $G$ 关于代数运算 $\circ$ 作成一个群。
\end{note}

\begin{theorem}[消去律与有限群判定]
设 $G$ 是一个乘法运算封闭的非空有限集合，如果 $G$ 满足结合律，且两个消去律成立，则 $G$ 是一个群。

\textbf{证：} 设 $G=\{a_1,a_2,\cdots,a_n\}$，对任意的 $a,b\in G$，考察 $aa_i$ 与 $aa_j$，如果 $aa_i=aa_j$，则由左消去律得，$a_i=a_j$，于是 $i=j$。这说明，$aa_1,aa_2,\cdots,aa_n$ 是 $G$ 中 $n$ 个不同的元素。

封闭有限集，有限有限不需要找单位元和逆元。消去律是未来解方程，群可以用消去律。
\end{theorem}

\begin{note}[消去律成立与证明G是一个群 重点]
封闭有限集 有限有限不需要找单位元和逆元 \\
消去律是未来解方程 \\
群可以用消去律
\end{note}

\begin{definition}[加群]
一个交换群叫做一个加群，是指我们把这个群的代数运算叫做加法，并且用符号 $+$ 来表示。
\end{definition}

% ===================== 第七节 =====================
\section{有限群与元素的阶}

\begin{definition}[有限群（定义 1.20）]
令 $(G,\cdot)$ 是一个群，我们称 $G$ 是一个有限群，当该群是有限的，则 $G$ 是有限的。

第一个概念是元素的阶，这与循环群密切相关，我们会慢慢引到循环群的概念。

令 $x\in G$，则 $x$（在 $G$ 中）的阶，记作 $|x|$，定义为那个最小的正整数 $n\in\mathbb{N}_1$，使得 $x^n=e$。若这样的 $n$ 不存在，则记 $|x|=\infty$。
\end{definition}

\begin{note}[群的阶与元素的阶]
群的阶是群中元素的个数。群中每一个元素的阶可以是不一样的。
\end{note}

\begin{note}[定义1 20元素的阶]
乘法不是真正的乘法 \\
$a$的$n$次方$=a\circ a\circ a$就是$n$次运算，运算可以定义 \\
阶是最小正整数 阶的符号和绝对值一样 \\
有限群每个元素的阶均有限 $a$和阶和$a^{-1}$的阶一样 \\
交换群可以拆括号
\end{note}

\begin{proposition}[命题 1.21]
$(G,\cdot)$ 是有限群，且 $x\in G$，则 $|x|<\infty$。换言之，有限群的每一个元素都可以通过自乘有限多次，都可以得到单位元。
\end{proposition}

\begin{note}[有限群每一个元素都有元]
不能不相等是因为如果都不相等就会出现无限了
\end{note}

\begin{proposition}[$a$ 和 $a^{-1}$ 阶数一样]
\[
|a|=|a^{-1}|
\]
若 $|a|=m$，那么 $a^m=e$，另外 $(a^{-1})^m=(a^m)^{-1}=e^{-1}=e$，知 $|a^{-1}|\le m$；

若 $|a^{-1}|=n\Rightarrow(a^{-1})^n=e\Rightarrow(a^n)^{-1}=e\Rightarrow a^n=e^{-1}=e$，$\therefore|a|\le n$，于是有 $n\le m$，

且 $m\le n\Rightarrow m=n$，即 $|a|=|a^{-1}|$。
\end{proposition}

\begin{note}[a和a的逆阶数一样]
设$a$的阶数是$m$，$a^{-1}$的$m$次方也是$e$，$m$大于等于$a^{-1}$的阶数 \\
同理设$a^{-1}$的阶数是$n$，$a$的$n$次方也是$e$，$n$大于等于$a$的阶数 \\
$m$大于等于$n$，$n$大于等于$m$
\end{note}

\begin{example}[阶的例题]
1. 若群 $G$ 的每一个元都适合方程 $x^2=e$，那么 $G$ 是交换群。

\textbf{解：} 令 $a$ 和 $b$ 是 $G$ 的任意两个元。由题设 $(ab)(ab)=(ab)^2=e$，另一方面 $(ab)(ba)=ab^2a=aea=a^2=e$，于是有 $(ab)(ab)=(ab)(ba)$。利用消去律，得 $ab=ba$，所以 $G$ 是交换群。
\end{example}

\begin{note}[命题1 2 6 带余除法和阶n的结合]
主要要把n拿出来
\end{note}

\begin{note}[阶数与奇偶]
由$a$和$a$的逆出现偶数，有$a$的逆是群的定义规定的
\end{note}

\begin{note}[阶是素数的群一定是循环]
也就是阶是素数的群每个不是单位元的都是生成元 先生成一个循环子群子群的阶是父群阶的因子 因为父群阶是素数所以子群阶就是父群的阶
\end{note}

% ===================== 第八节 =====================
\section{交换群（阿贝尔群）}

\begin{definition}[交换群]
若群 $G$ 的代数运算满足交换律，则称 $G$ 为交换群（Abel群）。
\end{definition}

\begin{example}
整数集 $Z$：

（1）对于数的普通加法：加法有结合律、交换律，左单位元 $0$，左逆元 $-a$，因此做成交换群，称为\textbf{整数加群}；

（2）对于数的普通乘法：左单位元 $1$，$a\ne\pm1$ 无逆元，不能做成群；

（3）对于运算 $a\circ b=a+b+4$：

$(a\circ b)\circ c=(a+b+4)\circ c=(a+b+4)+c+4=a+b+c+8$

$a\circ(b\circ c)=a\circ(b+c+4)=a+(b+c+4)+4=a+b+c+8$

有结合律，且有交换律。左单位元 $-4$：$-4\circ b=-4+b+4=b$；左逆元：$(-8-a)\circ a=-8-a+a+4=-4$。做成交换群。
\end{example}

\begin{definition}[定义 1.10：阿贝尔群]
若 $(G,\cdot)$ 是一个群，我们称它是阿贝尔群，或交换群，当该运算满足交换律，即
\[
\forall x,y\in G,\;x\cdot y=y\cdot x
\]
\end{definition}

\begin{note}[1 10阿贝尔群]
常见的群，加群最为普遍
\end{note}

\begin{note}[交换群 阿贝尔群]
交换群就是证明$ab=ba$
\end{note}

% ===================== 第九节 =====================
\section{子群}

\begin{definition}[子群]
一个群 $G$ 的一个子集 $H$ 叫做 $G$ 的一个子群，是指 $H$ 对于 $G$ 的乘法来说作成一个群。

设 $G$ 是群，$H$ 是 $G$ 的一个非空子集。如果按 $G$ 中元素的二元运算也构成群，则称 $H$ 是 $G$ 的一个子群，记为 $H\le G$。
\end{definition}

\begin{note}[子群]
继承运算 \\
子群也得是群而且运算和父空间一样
\end{note}

\begin{note}[定义子群的目的（定义 1.13）]
定义子群根本目的是证明它是一个群。
\end{note}

\begin{theorem}[定理 1：子群充要条件]
一个群 $G$ 的一个非空子集 $H$ 作成 $G$ 的一个子群的充要条件是：
\begin{enumerate}[(i)]
  \item $a,b\in H\Rightarrow ab\in H$；
  \item $a\in H\Rightarrow a^{-1}\in H$。
\end{enumerate}
\end{theorem}

\begin{note}[定理1子群充要条件]
子群单位元是原群的单位元，子群的逆元也是原群的逆元 \\
继承最重要的
\end{note}

\begin{theorem}[定理 2：子群充要条件]
一个群 $G$ 的一个非空子集 $H$ 作成 $G$ 的一个子群的充要条件是：
\begin{enumerate}[(iii)]
  \item $a,b\in H\Rightarrow ab^{-1}\in H$。
\end{enumerate}
\end{theorem}

\begin{note}[定理2非空子集充要条件]
定理2就是简化的定理1：把逆元条件和乘法封闭条件合并成一个条件 $ab^{-1}\in H$。
\end{note}

\begin{theorem}[定理 3：非空有限子群]
一个群 $G$ 的一个非空有限子集 $H$ 作成 $G$ 的一个子群的充要条件是
\[
a,b\in H\Rightarrow ab\in H.
\]
\end{theorem}

\begin{note}[学子群的意义]
\begin{enumerate}
  \item \textbf{理解群的结构：} 子群是群的"内部构件"。通过研究一个群的所有子群，我们可以了解群内部的分解和层次结构，这对分类和分析群的性质非常关键。例如，利用子群可以描述群的分解、直积和半直积结构，从而揭示群的内在对称性和代数性质。
  \item \textbf{简化证明和构造：} 在解决某些问题时，我们可以先在一个子群内证明结论，然后再推广到整个群。例如，在证明群的某些性质时，先考虑生成的循环子群或正规子群往往会使证明过程简单明了。
  \item \textbf{群的同态与商群：} 子群，尤其是正规子群，是构造群同态和商群的基础。通过研究子群，我们可以定义群的同态映射、了解核与像的关系，从而构建新的群（商群），这对于群的分类、群论的应用及现代代数的发展都非常重要。
\end{enumerate}
\end{note}

\begin{proposition}[命题 1.11]
有时为了方便起见，我们把子群的后两个条件合并，缩成两个条件，即
\[
e\in H,\qquad\forall x,y\in H,\;x\cdot y^{-1}\in H
\]
这条命题说的便是，这是子群的等价条件。
\end{proposition}

\begin{note}[命题1 11证明]
第二个条件x可以是e，y则是H中任意元素证明了每一个都有逆元再证明x.y属于H即可
\end{note}

\begin{note}[子群里面的元素当陪集的x，则这个陪集还是原来子群xH=H]
这里就把前面的g可以分为H里面的和G-H里面的，核心原因是子群对乘法封闭
\end{note}

\begin{note}[命题1 12特殊线性群是群]
因为它的定义就是行列式是1，$1\times1$永远是1所以封闭性极好逆元也是自己啥都有了
\end{note}

% ===================== 第十节 =====================
\section{陪集与拉格朗日定理}

\begin{definition}[左陪集（定义 1.27）]
令 $G$ 是一个群，$H<G$ 是一个子群，$a\in G$。则 $aH$ 是 $H$ 的一个左陪集（由 $a$ 引出），定义为
\[
aH=\{ax:x\in H\}
\]
也就是说，由 $a$ 引出的 $H$ 的左陪集，就是用 $a$ 左乘了 $H$ 中的每一个元素所得到的。特别地，$eH=H$ 也是一个左陪集。
\end{definition}

\begin{definition}[右陪集]
一个群 $G$ 和 $G$ 的一个子群 $H$。我们规定一个 $G$ 的元中间的关系 $\sim$：$a\sim b$，当而且只当 $ab^{-1}\in H$ 的时候。
\begin{enumerate}
  \item $aa^{-1}=e\in H$，所以 $a\sim a$；
  \item $ab^{-1}\in H\Rightarrow(ab^{-1})^{-1}=ba^{-1}\in H$，所以 $a\sim b\Rightarrow b\sim a$；
  \item $ab^{-1}\in H$，$bc^{-1}\in H\Rightarrow(ab^{-1})(bc^{-1})=ac^{-1}\in H$，所以 $a\sim b$，$b\sim c\Rightarrow a\sim c$。
\end{enumerate}
由上面的等价关系 $\sim$ 所决定的类叫做子群 $H$ 的\textbf{右陪集}。包含元 $a$ 的右陪集用符号 $Ha$ 来表示。
\[
H\le G,\;a\in G\text{ 则 }Ha=\{ha\mid h\in H\}
\]
\end{definition}

\begin{note}[右陪集的定义]
陪集是陪伴着出现的没有对象就不存在陪伴 \\
$a$和$b$满足这个关系当且仅当$ab^{-1}$在$H$子群中 \\
下面的三条是证明这个关系是等价关系 \\
$Ha$就是$H$中每一个元素乘$a$；$aH$一般不是群但是可以与$H$构成双射
\end{note}

\begin{lemma}[引理 1.2]
令 $G$ 是一个有限群，$H<G$ 是一个子群，$a\in G$。我们通过左乘 $a$ 来定义 $f:H\to aH$，$f(x)=ax$，则 $f$ 是一个双射。特别地，$|H|=|aH|$。
\end{lemma}

\begin{proposition}[命题 1.33：陪集要不然相等要不然无交]
令 $G$ 是一个有限群，$H<G$ 是一个子群，$a,b\in G$。则左陪集 $aH$ 和 $bH$ 要么相等，要么无交。也就是说，我们有 $aH=bH$，或 $aH\cap bH=\varnothing$。
\end{proposition}

\begin{note}[命题1 3 3子群陪集要不然相等要不然无交]
$a_1=bh_2\cdot h_1^{-1}$ 两边都能加h \\
比如说剩余类加群$\{0,1,2,3,4,5\}$子群，有1阶的就是$\{0\}$，2阶$\{0,3\}$生成元是3
\end{note}

\begin{example}
例子：群 $\mathbb{Z}_4$ 与子群 $H=\{0,2\}$

考虑加法模 4 的群 $\mathbb{Z}_4=\{0,1,2,3\}$，其运算为加法模 4。子群 $H=\{0,2\}$ 是 $\mathbb{Z}_4$ 的一个子群。

$0+H=\{0+0,0+2\}=\{0,2\}$；$2+H=\{2+0,2+2\}=\{2,4\}\bmod 4=\{2,0\}=\{0,2\}$。

可以看到，$0+H=2+H=\{0,2\}$，这说明两个不同的元素 0 和 2 所生成的左陪集是相等的。
\end{example}

\begin{note}[子群里面的元素当陪集的 $x$，则这个陪集还是原来子群 $xH=H$]
子群里面的元素当陪集的 $x$，则这个陪集还是原来子群 $xH=H$。
\end{note}

\begin{proposition}[命题 1.34、1.35：嵌套子群的陪集]
\textbf{命题 1.34：} 令 $K<H<G$ 是三个有限群，则
\[
[G:K]=[G:H][H:K]
\]

\textbf{命题 1.35：} 令 $(G,\cdot)$ 是一个群，若 $H,K<G$ 是两个有限子群，则
\[
|HK|=\frac{|H||K|}{|H\cap K|}
\]
其中 $HK=\{hk:h\in H,k\in K\}$。$<$ 是子群的意思。
\end{proposition}

\begin{note}[命题1 34 1 35嵌套子群的陪集]
$<$是子群的意思 \\
证明1 $|G|/|H|\cdot|H|/|E|$ \\
证明2 I J表示的就是能让他们形成不相交陪集的代表元，需要证明两件事一件事是枚举了所有陪集也就能覆盖G，第二件事就是证明集合两两不交
\end{note}

\begin{definition}[商集与指数（定义 1.28）]
一个群 $G$ 的一个子群 $H$ 的右陪集（或左陪集）的个数叫做 $H$ 在 $G$ 里的\textbf{指数}。

下面我们要用右陪集来证明几个重要定理。因为左陪集和右陪集的对称性，凡是我们以下用右陪集的地方也都可以用左陪集来代替。

引理：一个子群 $H$ 与 $H$ 的每一个右陪集 $Ha$ 之间都存在一个一一映射。证明：$\varphi:h\to ha$ 是 $H$ 与 $Ha$ 间的一一映射。
\end{definition}

\begin{note}[定义1 28商集与指数]
指数就是左陪集个数
\end{note}

\begin{theorem}[定理 1：左右陪集个数相等]
一个子群 $H$ 的右陪集的个数和左陪集的个数相等：它们或者都是无限，或者都有限并且相等。
\end{theorem}

\begin{note}[定理1 左右陪集个数相等]
先构造左陪集和右陪集的一个映射 \\
i是证明a，b的选取不影响
\end{note}

\begin{theorem}[拉格朗日定理（定理 1.1）]
令 $G$ 是一个有限群，$H<G$ 是一个子群，则
\[
|G|=[G:H]|H|
\]
特别地，$|H|\,\big|\,|G|$。
\end{theorem}

\begin{note}[拉格朗日定理]
$G$的阶等于商群的数量乘子群的阶 \\
比如循环群$\{0,1,2,3,4,5\}$ 子群$\{0,3\}$ 他生成的商群$\{1,4\}$ $\{2,5\}$ $\{0,3\}$ $6=3\times2$ \\
子群的阶一定是父群阶的因数
\end{note}

% ===================== 第十一节 =====================
\section{正规子群与商群}

\begin{definition}[正规子群（不变子群）]
一个群 $G$ 的一个子群 $N$ 叫做一个\textbf{正规子群}（不变子群），是指对于 $G$ 的每一个元 $a$ 来说，都有
\[
Na=aN.
\]
一个正规子群 $N$ 的一个左（或右）陪集叫做 $N$ 的一个\textbf{陪集}。$Na=aN$ 是用于验证的。

任何两个群的两个平凡子群都是不变子群；交换群的子群都是不变子群；素数阶的群都是循环群。代表选取不影响是重点。
\end{definition}

\begin{definition}[正规子群（定义 1.29）]
$N\triangleleft G$。
\end{definition}

\begin{note}[定义1 29正规子群]
得给商集 群结构，$a$和$a'$是不同代表元，正规子群才是三角，普通子群只是小于等于
\end{note}

\begin{theorem}[定理 1：不变子群的判定]
一个群 $G$ 的一个子群 $N$ 是一个不变子群的充分而且必要条件是：
\[
aNa^{-1}=N
\]
对于 $G$ 的任意一个元 $a$ 都对。
\end{theorem}

\begin{theorem}[定理 2：不变子群判定方式 2]
一个群 $G$ 的一个子群 $N$ 是一个不变子群的充分而且必要条件是：
\[
a\in G,\;n\in N\Rightarrow ana^{-1}\in N
\]
\end{theorem}

\begin{lemma}[引理 1.4：正规子群]
正规子群的意义就是为了让陪集成为群（商群）。$N\trianglelefteq G$ 使得陪集集合 $G/N$ 上的乘法良定义，从而构成商群。
\end{lemma}

\begin{note}[正规子群定义]
任何两个群的两个平凡子群都是不变子群 \\
交换群的子群都是不变子群 \\
素数阶的群都是循环群 \\
代表选取不影响是重点
\end{note}

\begin{note}[正规子群的意义就是为了让看看陪集是不是群]
加群逆元就是相反数，乘群逆元就是倒数
\end{note}

\begin{note}[指数是2的子群一定是不变的]
也就是要不然g就在eN里面左乘右乘没区别，要不然因为已经有一个eN了右陪集，左陪集只有一个一个座位所以他们只能是一个人
\end{note}

\begin{proposition}[命题 1.36：良定义]
令 $(G,\cdot)$ 是一个群，且 $N\trianglelefteq G$，$a,b\in G$，则
\[
(aN)\cdot(bN)=(ab)N
\]
是良定义的。
\end{proposition}

\begin{note}[命题1 36良定义]
正规子群就是良定义，良定义在数学里面就是不同表达方式下的取值说一样的，和选取的代表元无关
\end{note}

\begin{proposition}[命题 1.37：商群]
令 $(G,\cdot)$ 是一个群，且 $N\trianglelefteq G$，则 $(G/N,\cdot)$ 构成一个群，称为（$G$ 在 $N$ 上的）商群，其中单位元是 $eN=N$，每个陪集 $aN$ 的逆元是 $a^{-1}N$。
\end{proposition}

\begin{theorem}[定理 3：不变子群的陪集叫商群]
一个不变子群的陪集对于上边规定的乘法来说作成一个群。

定义：一个群 $G$ 的一个不变子群 $N$ 的陪集所作成的群叫做一个商群。这个群我们用符号 $G/N$ 来表示。
\end{theorem}

\begin{proposition}[命题 1.38：正规子群的交集还是正规子群]
令 $\{N_i\}_{i\in I}$ 是一族 $G$ 的正规子群，则它们的交集仍然是 $G$ 的正规子群，即
\[
\bigcap_{i\in I}N_i\trianglelefteq G
\]
\end{proposition}

\begin{proposition}[命题 1.39、1.40：$e$ 和自己都是自己的正规子群；交换群的子群都是正规子群]
$e$ 和自己都是自己的正规子群；交换群的子群都是正规子群。
\end{proposition}

\begin{example}[不变子群的中心]
$N$ 刚好包含群 $G$ 的所有有以下性质的元 $n$：$na=an$，不管 $a$ 是 $G$ 的哪一个元，那么 $N$ 是 $G$ 的一个不变子群。这个不变子群叫做 $G$ 的\textbf{中心}。
\end{example}

\begin{note}[商群性质]
商群有下列常用的性质：
\begin{enumerate}
  \item 商群 $G/N$ 的阶 $=[G:N]$；
  \item 如果 $G$ 是有限群，则商群 $G/N$ 的阶 $=[G:N]=\dfrac{|G|}{|N|}$；
  \item 有限群的商群还是有限群，且其任一商群的阶是群阶数的因数；
  \item $N\trianglelefteq G$，则 $\bar{e}=eN=N$ 为商群 $G/N$ 的单位元，$a^{-1}N$ 为 $aN$ 的逆元。
\end{enumerate}
\end{note}

% ===================== 第十二节 =====================
\section{同态与同构}

\begin{definition}[同态与同构]
$\phi$ 是 $A$ 与 $\bar{A}$ 间的一一映射，$\circ$ 和 $\bar{\circ}$ 分别是 $A$ 和 $\bar{A}$ 的代数运算，如果 $\forall a,b\in A$，只要 $a\to\bar{a}$，$b\to\bar{b}$，就有 $a\circ b\to\bar{a}\bar{\circ}\bar{b}$，则说 $\phi$ 是 $A$ 与 $\bar{A}$ 间的同构映射，并说 $A$ 与 $\bar{A}$ 同构，记为 $A\cong\bar{A}$。

同构只关心运算规律，不关心元素长什么样，可以把复杂的群换成结构相同但更容易处理的群。两个集合中的代数运算可以不一样，但是否能够构造群同态或同构，取决于它们的代数结构如何关联。
\end{definition}

\begin{example}[同态例子]
例1. 设 $A=\{\text{所有实数 }x\}$，$A$ 的代数运算是实数的乘法。以下映射是否是同态？

（1）若 $f:x\to|x|$。（1）$f$ 是同态。

（2）若 $f:x\to2x$。（2）$f$ 不是同态。
\end{example}

\begin{theorem}[同态定理]
两个代数系统 $G$ 与 $\bar{G}$ 同态，若 $G$ 是群，则 $\bar{G}$ 也是群。
\end{theorem}

\begin{note}[同态的意义]
同构：同构只关心运算规律，不关心元素长什么样。它能把"复杂的群"换成"结构相同但更容易处理的群"。
\end{note}

\begin{note}[同态同构的意义]
从这里，我们有两个延伸。第一个延伸是，什么是"引出"？我们可以从一个元素引出一个同态，实际上，也可以从任何一个元素引出一个原幺半群的子幺半群。另一方面，其实每一个子集都可以引出一个子幺半群。第二个延伸是"同构"。同构首先是一个双射，这意味着你可以在元素之间建立一一对应；在双射的基础上，同构的含义是，换了标签下，运算是不变的——也就是说，这两个结构（比如幺半群）的唯一区别就是"标签不同"。
\end{note}

\begin{definition}[群同态（定义 1.14）]
令 $(G,\cdot)$，$(G',*)$ 是两个群，且 $f:G\to G'$ 是一个映射，我们称 $f$ 是一个群同态，当其保持了乘法运算，即
\[
\forall x,y\in G,\;f(x\cdot y)=f(x)*f(y)
\]
\end{definition}

\begin{note}[定义1 16群同态]
左乘，同态最好的就是可以拆括号让x和$x^{-1}$相遇
\end{note}

\begin{proposition}[命题 1.13：群同态单位元对单位元，逆元对逆元]
若 $f:(G,\cdot)\to(G',*)$ 是一个群同态，则 $f(e)=e'$，$f(x^{-1})=f(x)^{-1}$。
\end{proposition}

\begin{note}[命题1 13群同态单位元对单位元逆元对逆元]
证明思路因为$e\cdot e=e$所以要添项减少项目的时候e就是很有用e也可以拆成两个互逆的元素相乘把我们想要的东西引入
\end{note}

\begin{proposition}[命题 1.14：$\det:GL(n,\mathbb{R})\to(\mathbb{R}^\times,\cdot)$ 是一个乘法群同态]
$\det:GL(n,\mathbb{R})\to(\mathbb{R}^\times,\cdot)$ 是一个乘法群同态。
\end{proposition}

\begin{definition}[核与像（定义 1.15）]
群同态下定义：核是映射到 $e'$ 的原像集，$\operatorname{im}f$ 是像集。
\[
\ker(f)=\{x\in G\mid f(x)=e'\},\qquad\operatorname{im}(f)=\{f(x)\mid x\in G\}
\]
\end{definition}

\begin{note}[定义1 15群同态下定义 核是映射到e'的原像集 imf是像集]
核是映射到对方单位元对应的原像集，像是G能映射到$G'$的所有$G'$的元素。但是因为同态不要求单满所以$\operatorname{im}(f)$至少G的子集 \\
要证两件事1有单位元2.$x,y$在这里面那么$xy^{-1}$次方也在，用$\ker f$的定义也就是验证这个东西能不能映射到$e'$上。$e'$本身就是一个群，反正要把运算通过同态的性质转化到已知的群上
\end{note}

\begin{proposition}[命题 1.16]
令 $f:(G,\cdot)\to(G',*)$ 是一个群同态，则 $f$ 是一个单同态当且仅当 $\ker(f)=\{e\}$。也就是说，一个群同态是单的当且仅当核是平凡的。
\end{proposition}

\begin{proposition}[命题 1.17：同态就是拆括号的游戏]
若 $f:(G,\cdot)\to(G,*)$ 是一个群同构，则 $f^{-1}$ 也是群同构。
\end{proposition}

\begin{note}[命题1 17同态就是拆括号的游戏]
同态就是拆括号合括号的游戏
\end{note}

\begin{note}[同态的定义]
同态映射 A这边对两个元素操作A得到的 和$A'$中对这两个元素在$A'$中的映射操作A 得到的还是映射 \\
同态映射+满射才能说A与$A'$是同态 \\
同态传递结合律分配律交换律
\end{note}

\begin{note}[同态同构的意义]
同态保持运算，同构完全一致
\end{note}

\begin{note}[同构]
同构=同态+限制必须一一对应既单又满 \\
同构只看运算规律，把复杂集合的运算同构简化 \\
自己到自己叫自同构
\end{note}

\begin{note}[构造同态同构的关键点]
\textbf{同态（Homomorphism）：} 只需要映射保持运算结构：$\phi(a\cdot b)=\phi(a)*\phi(b)$。例如，$(\mathbb{R},+)$ 和 $(\mathbb{R}^+,\times)$ 的运算不同，但可以用指数映射 $\phi(x)=e^x$ 连接它们，使其成为同构。

\textbf{同构（Isomorphism）：} 除了运算关系必须保持，映射还必须是双射。同构要求双射，且必须使运算方式保持一致，即使运算符号不同，也可以通过映射来建立联系。

证明同构：证明任意 $\phi(a\circ b)=\phi(c)$，同时 $\phi(a)\circ\phi(b)=d$；证明单射反证就是一对多；特殊点就是切入口。
\end{note}

% ===================== 第十三节 =====================
\section{群同构定理}

\begin{theorem}[群同构第一定理（定理 1.2）]
令 $f:G\to G'$ 是一个群同态，则 $\ker(f)\triangleleft G$，且 $G$ 在 $\ker(f)$ 上的商群同构于 $\operatorname{im}(f)$，即
\[
G/\ker(f)\simeq\operatorname{im}(f)
\]
特别地，若 $f$ 是满同态，则
\[
G/\ker(f)\simeq G'
\]
若 $f$ 是单同态，则
\[
G/\{e\}\simeq G\simeq\operatorname{im}(f)
\]
若 $G$ 是有限群，则
\[
\frac{|G|}{|\ker(f)|}=|\operatorname{im}(f)|
\]
\end{theorem}

\begin{note}[群同构第一定理1 2]
$\ker(f)$是映射到e的原像集，G上由$\ker(f)$的陪集组成群和$\operatorname{im}(f)$也就是值域群上同构 \\
如果f是满射，$\operatorname{im}(f)$就是$G'$，如果f是单同态$\ker(f)=\{e\}$ \\
如果G有限则他的阶等于$\ker(f)$的阶乘$\operatorname{im}(f)$的乘积
\end{note}

\begin{theorem}[群同构第二定理（定理 1.3）]
令 $(G,\cdot)$ 是一个群，且 $N\trianglelefteq G$，$H<G$。则 $H\cap N\trianglelefteq H$，$N\trianglelefteq HN$，且
\[
H/(H\cap N)\simeq HN/N
\]
这和之前两个子群乘积的阶的公式是类似的。
\end{theorem}

\begin{note}[群同构第二定理]
N是正规子群，H是真子群 H和N的交集是H的正规子群 \\
$H\cap N$构成的子群的陪集组成的商群和H中元素与N组合构成的陪集的商群是同构的
\end{note}

\begin{note}[证明]
正规子群一共干的就是一件事从父群里面任意挑出a和对应的a的逆，再证明正规子群里面的元素可以达成$aNa^{-1}=N$，而证明的时候就要用增加a和$a^{-1}$把想要的形式凑出来
\end{note}

\begin{note}[证明]
$\ker(f)$被拿出来就是因为它能映射到$e'$只要映射到$G'$这一项就会变成$e'$然后消失掉
\end{note}

\begin{theorem}[群同构第三定理（定理 1.4）]
令 $(G,\cdot)$ 是一个群，且 $N\trianglelefteq G$，$M\trianglelefteq G$，$M<N$。则 $N/M\trianglelefteq G/M$，且
\[
(G/M)/(N/M)\simeq G/N
\]
\end{theorem}

% ===================== 第十四节 =====================
\section{变换群与置换群}

\begin{definition}[变换群]
\textbf{1 变换：} 设 $M$ 是一个非空集合，若 $\tau$ 是 $M$ 到 $M$ 上的映射 $\tau:M\to M$，就称 $\tau$ 是 $M$ 的一个\textbf{变换}。

\textbf{2 变换乘法：} $\tau_1$，$\tau_2$ 是 $M$ 的变换，规定 $\forall a\in M$：
\[
\tau_1\tau_2:a\to(a^{\tau_1})^{\tau_2}=a^{\tau_1\tau_2}
\]
称 $\tau_1\tau_2$ 为 $\tau_1,\tau_2$ 的乘法。

\textbf{3 恒等变换：} $\varepsilon:\forall\tau,\;\varepsilon\tau=\tau\varepsilon=\tau$。

\textbf{定理 1：} 假定 $G$ 是集合 $A$ 的若干个变换所作成的集合，并且 $G$ 包含恒等变换 $\varepsilon$。若是对于上述乘法来说 $G$ 作成一个群，那么 $G$ 只包含 $A$ 的一一变换。

变换得自己往自己身上映射；变换群里面的元素是变换。
\end{definition}

\begin{note}[变换群]
变换得自己往自己身上映射 \\
变换乘法 恒等变换 \\
变换群里面的元素是变换
\end{note}

\begin{note}[变换群证明]
（1）非空、代数运算、结合律都满足，有单位元 $\varepsilon$。那么"逆元"问题能解决吗？事实上，$\tau_1$ 就没有逆元。因为如果 $\tau_1$ 有逆元 $\tau$，那么必有 $\tau_1\tau=\varepsilon$ 且 $\tau\tau_1=\varepsilon$。但是
\[
(1)^\varepsilon=(1)^{\tau_1\tau}=(1^{\tau_1})^\tau=(1)^\tau\Rightarrow(1)^\tau=1
\]
而
\[
(2)^\varepsilon=(2)^{\tau_1\tau}=(2^{\tau_1})^\tau=(1)^\tau\Rightarrow(1)^\tau=2
\]
导致矛盾，故 $\tau_1$ 没有逆元。因此 $T(M)$ 不能成为群。
\end{note}

\begin{note}[变化群的意义]
可以通过一个简单的例子来理解。设想你有一个正方形，你可以对这个正方形进行各种操作（旋转、翻转等），这些操作我们称为"变换"。变换的组合（封闭性）：如果你先对正方形进行一次变换，再进行另一次变换，整体效果依然是一个变换。比如先旋转90°再旋转90°，结果等于旋转180°。
\end{note}

\begin{definition}[置换与置换群]
一个有限集的一个一一变换叫做一个置换。一个有限集的若干个置换作成的一个群叫做一个置换群。

我们看一个有限集合 $A$，$A$ 有 $n$ 个元 $a_1,a_2,\cdots,a_n$。由 §5 定理 2，$A$ 的全体置换作成一个群 $G$。
\end{definition}

\begin{note}[置换与置换群]
有限集+变换=置换 \\
一一变换为置换
\end{note}

\begin{definition}[置换的书写方式]
置换
\[
\pi:a_i\longrightarrow a_{k_i},\quad i=1,2,\cdots,n
\]
这样一个置换所发生的作用完全可以由 $(1,k_1),(2,k_2),\cdots,(n,k_n)$ 这 $n$ 对整数来决定。我们表示重量的第一个方法就是把这个置换写成
\[
\begin{pmatrix}1&2&\cdots&n\\k_1&k_2&\cdots&k_n\end{pmatrix}
\]
例1 $n=3$，$\pi:1\to2,\;2\to3,\;3\to1$，那么
\[
\pi=\begin{pmatrix}1&2&3\\2&3&1\end{pmatrix}
\]
\end{definition}

\begin{note}[循环置换]
3循环 第k个接回第一个 \\
不变的映射不写 \\
两个循环不能有公共的东西加不相连的循环
\end{note}

\begin{note}[置换群]
置换是一一对应的变化，不相连就是交集是空，写成循环置换就是为了省地方 \\
n次对称群就是全体置换
\end{note}

\begin{note}[只要不相连就可以快速求阶数而且有交换律]
$ab=ba$，$(123)(45)$，$3\times2=6$。
\end{note}

\begin{definition}[对称群]
一个包含 $n$ 个元的集合的全体置换作成的群叫做 $n$ 次对称群，这个群我们用 $S_n$ 来表示。
\end{definition}

\begin{theorem}[$n$ 次对称群 $S_n$ 的阶是 $n!$]
定理 1：$n$ 次对称群 $S_n$ 的阶是 $n!$。

现在我们要看一看表示一个置换的符号。$n$ 个元的置换一共有 $n!$ 个。由初等代数我们知道，$n$ 个元的置换一共有 $n!$ 个，要作 $n$ 个元的一个置换，就是要替每一个元选定一个对象，我们替 $a_1$ 选定对象时，有 $n$ 个可能，选定了以后，再替 $a_2$ 选时，就只有 $n-1$ 个可能，这样下去，一共可以得到 $n(n-1)\cdots2\cdot1=n!$ 个不同的置换。
\end{theorem}

% ===================== 第十五节 =====================
\section{凯莱定理}

\begin{theorem}[凯莱定理：任何群都是同一个变换群的同构]
定理（凯莱定理）任何群都同一个变换群同构。

\textbf{证明：} 假定 $G$ 是一个群，$G$ 的元是 $a,b,c,\cdots$。我们在 $G$ 里任意取出一个元 $x$ 来，那么
\[
g\longrightarrow gx=g^{\tau_x}
\]
是集合 $G$ 的一个变换，因为给了 $G$ 的任意元 $g$，我们能够得到一个唯一的 $G$ 的元 $g^{\tau_x}$，这样当 $G$ 的每个元 $x$ 可以得到一个 $G$ 的变换放在一起，作成一个集合 $\tilde{G}=\{\tau_a,\tau_b,\tau_c,\cdots\}$。那么
\[
\phi:\quad x\longrightarrow\tau_x
\]
是 $G$ 到 $\tilde{G}$ 间的满射，且消去律告诉我们：$x\ne y\Rightarrow gx\ne gy$ 告诉我们，$\tau_x\ne\tau_y$。所以 $\phi$ 是 $G$ 与 $\tilde{G}$ 间的一一映射。
\end{theorem}

\begin{note}[凯莱定理任何群都是同一个群的同构]
第一步构造映射g到gx这个操作叫做$\tau_x$ \\
每一个元都有一种变换 再构建变换和元素之间的映射 \\
gx之间是由定义的运算连接的而不是乘法 \\
这步证明的是$g(x\circ y)=g(x)\circ g(y)$
\end{note}

% ===================== 第十六节 =====================
\section{循环群}

\begin{definition}[定义 1.22：$x$ 的幂]
令 $(G,\cdot)$ 是一个群，且 $x\in G$。若 $n\in\mathbb{N}_1$，我们定义 $x^{-n}=(x^{-1})^n$，另外 $x^0=e$。

命题 1.22：令 $(G,\cdot)$ 是一个群，任取 $x\in G$，则 $f:(\mathbb{Z},+)\to(G,\cdot)$，定义为 $f(n)=x^n$，是一个群同态。

$x$ 的幂是群同态，$f$ 是同态映射。
\end{definition}

\begin{definition}[由 $x$ 生成子群（定义 1.23）]
令 $(G,\cdot)$ 是一个群，且 $x\in G$，则 $\langle x\rangle$，被称为由 $x$ 生成的群，定义为
\[
\langle x\rangle=\{x^n:n\in\mathbb{Z}\}
\]
利用同态的性质，我们有的直接推论就是，刚才提到的同态的像，也就是 $\{x^n:n\in\mathbb{Z}\}$ 是一个子群，被称为由 $x$ 生成的群。

阶只是 $n$ 次方可以等于 $e$，没有说 $n$ 次方的过程，这 $n$ 个元素就是群中所有的元素。
\end{definition}

\begin{definition}[循环群（定义 1.25）]
令 $(G,\cdot)$ 是一个群。若存在 $x\in G$，使得 $G=\langle x\rangle$，则 $G$ 被称为一个循环群，而 $x$ 被称为 $G$ 的一个生成元。

循环群的关键是找到固定元，由固定元可以生成所有的元素。若群 $G$ 中每个元都能表示成某个固定元 $a$ 的乘方，就称群 $G$ 为循环群，也称群 $G$ 为由元 $a$ 生成的群，记为 $G=(a)$，称 $a$ 是 $G$ 的一个生成元。

\textbf{例1：} 整数加群 $Z$ 是无限阶循环群。$\because 1\in Z$，$m=\underbrace{1+1+\cdots+1}_{m\text{个}}=1^m$，$-m=\underbrace{(-1)+(-1)+\cdots+(-1)}_{m\text{个}}=(-1)^m=1^{-m}$，$0=1^0$，$\therefore Z=(1)$。

\textbf{例2：} 模 $n$ 的剩余类加群 $Z_n=([1])$ 是 $n$ 阶循环群。
\end{definition}

\begin{note}[循环群 1]
循环群的关键是找到固定元，由固定元可以生成所有的元素，m是加法做多少次也就是运算多少次 \\
余k的可由余1的自+k次
\end{note}

\begin{note}[循环群 2]
由一个元形成群所有元素 所有元素都可以写成生成元的次方
\end{note}

\begin{proposition}[命题 1.26：有限循环群的结构]
令 $G=\langle x\rangle$ 是有限循环群，假设 $|x|=n$，则 $G=\{e,x,x^2,\cdots,x^{n-1}\}$，其中乘法法中的这些元素是两两不同的。
\end{proposition}

\begin{note}[循环子群]
循环群的阶是父群阶的因数，阶也就是个数
\end{note}

\begin{proposition}[命题 1.30：循环群的生成元]
令 $G=\langle x\rangle$ 是一个 $n$ 阶循环群。假设 $1\le m\le n$，则 $x^m$ 的阶为
\[
|x^m|=\frac{n}{\gcd(n,m)}
\]
\textbf{命题 1.31：} 令 $G=\langle x\rangle$ 是一个 $n$ 阶循环群，则 $x^m$（$1\le m\le n$）是个生成元，当且仅当 $\gcd(m,n)=1$。根据欧拉 $\phi$ 函数的定义，这些生成元的个数正是 $\phi(n)$。所以循环群的生成元不止一个。
\end{proposition}

\begin{note}[命题1 30循环群的生成元的幂指数和循环群阶数互素]
所以循环群的生成元不止一个
\end{note}

\begin{proposition}[$n$ 阶循环群互相同构，无限循环群彼此同构（命题 1.2.9）]
令 $G=\langle x\rangle$ 是有限循环群，$|x|=n$，则 $G\cong\mathbb{Z}/n\mathbb{Z}$。无限循环群 $G=\langle x\rangle$（$|x|=\infty$）则 $G\cong\mathbb{Z}$。故所有 $n$ 阶循环群互相同构，所有无限循环群彼此同构。
\end{proposition}

\begin{note}[循环群的同构]
把生成元做映射 \\
同构证明是映射证明同态+单+满
\end{note}

\begin{note}[循环群的子群还是循环群]
其实也就是要证明子群中最小的正幂次方$a^k$是生成元，设其他的元素不整除则$m=kq+r$群有逆元 所以$a^r$也在循环群内 因为r比k小，k又是最小的正幂r只能等于0所有元都是k的整数倍
\end{note}

\begin{note}[无限阶的群和任意循环群同态]
无限阶循环群（如$\mathbb{Z}$）可以同态到任意循环群：同态只需要满射，不必是单射，多对一可以，一对多不行。
\end{note}

% ===================== 第十七节 =====================
\section{直积}

\begin{definition}[两个群直积（定义 1.18）]
令 $(G,\cdot_1)$，$(G',\cdot_2)$ 是两个群，我们构造它们的直积，不妨记 $(G\times G',*)$。对于 $(x,y),(x',y')\in G\times G'$，我们定义逐坐标的乘积，为
\[
(x,y)*(x',y')=(x\cdot_1 x',\;y\cdot_2 y')
\]

\textbf{命题 1.18：} 若 $(G,\cdot_1)$，$(G',\cdot_2)$ 是两个群，则它们的直积 $(G\times G',*)$ 还是一个群。

\textbf{证明：} 封闭性：因为 $G$ 在 $\cdot_1$ 下封闭，$G'$ 在 $\cdot_2$ 下封闭，而 $G\times G'$ 的元素乘积是逐坐标定义的，则 $G\times G'$ 在 $*=(\cdot_1,\cdot_2)$ 下也是封闭的。结合律：同样，逐坐标有结合律，故整体也有结合律。单位元：$(e,e')$ 是直积的单位元。逆元：对于任意 $(x,y)\in G\times G'$，$(x^{-1},y^{-1})$ 是 $(x,y)$ 的逆元。
\end{definition}

\begin{definition}[一族群的直积（定义 1.19）]
令 $(G_i,\cdot_i)$ 是一族群，其中 $I$ 是一个指标集，我们定义它们的直积为 $(\prod_{i\in I}G_i,*)$，同样通过逐点的乘积，对于 $(x_i)_{i\in I},(y_i)_{i\in I}\in\prod_{i\in I}G_i$，我们定义
\[
(x_i)_{i\in I}*(y_i)_{i\in I}=(x_i\cdot_i y_i)_{i\in I}
\]
\textbf{命题 1.19：} 若 $(G_i,\cdot_i)$ 是一族群，则它们的直积 $(\prod_{i\in I}G_i,*)$ 还是一个群。

既然有了两个群的直积，我们就可以类似地定义有限多个群的直积，同样通过逐坐标的乘法。甚至进一步，我们可以对于任意一族群，定义它们的直积。
\end{definition}

\begin{definition}[$n$ 阶一般线性群（定义 1.11）]
我们对于那些 $n*n$ 可逆实矩阵构成的乘法群，称为（实数上的）$n$ 阶一般线性群，记作 $(GL(n,\mathbb{R}),\cdot)$。由于一个矩阵可逆当且仅当其行列式不为零，因此
\[
GL(n,\mathbb{R})=\{A\in M(n,\mathbb{R}):\det(A)\ne 0\}
\]
注意，这的确是个群，因为本身 $n*n$ 实矩阵对乘法构成幺半群（单位元是单位矩阵），接下来我们取所有可逆元素，就得到了一般线性群。提到了一般线性群很难不去提特殊线性群。
\end{definition}

\begin{definition}[$n$ 阶特殊线性群行列式是 1（定义 1.12）]
（实数上的）$n$ 阶特殊线性群是由那些行列式恰好是 $1$ 的 $n*n$ 实矩阵构成的乘法群，记作 $(SL(n,\mathbb{R}),\cdot)$，即
\[
SL(n,\mathbb{R})=\{A\in M(n,\mathbb{R}):\det(A)=1\}
\]
注意，这个定义是尚且不良好的，因为我们并没有证明特殊线性群是个群。实际上，它是个子群。我们为什么不给出子群的条件呢？正如我们给出子幺半群的条件那样。
\end{definition}

\begin{proposition}[命题 1.12：特殊线性群是群]
$SL(n,\mathbb{R})$ 是 $GL(n,\mathbb{R})$ 的一个子群，从而是一个群。（由行列式同态 $\det:GL(n,\mathbb{R})\to(\mathbb{R}^\times,\cdot)$ 的核即得。）
\end{proposition}

\begin{example}[4 元群]
假定 $F$ 是一个有四个元的域。证明：(i) $F$ 的特征是 $2$；(ii) $F$ 的不等于 $0$ 或 $1$ 的两个元都适合方程 $x^2=x+1$。

\textbf{解：} (a) $F$ 的特征是 $F$ 的非零元（对加法来说的）相同的阶，并且是一个素数。$F$ 作为加群的阶是 $4$，因此 $F$ 的非零元（对加法来说）的阶只能是 $1$，$2$ 或 $4$，其中只有 $2$ 是素数，所以特征是 $2$。因此加群 $F$ 与克莱因四元群 $B_4$ 同构。

(b) 另一方面乘群 $F^\times$ 的阶是 $3$，因而是一个循环群 $(a)$，而 $F^\times$ 的元是 $1,a,a^2$。这样 $F=\{0,1,a,a^2\}$。由于加群 $F$ 与 $B_4$ 同构，有 $a+1=a^2$，$a^2+1=a=(a^2)^2$，因此 $F$ 的不等于 $0$ 或 $1$ 的两个元 $a$ 和 $a^2$ 都适合方程 $x^2=x+1$。有限域一定有阶，阶只能是 $4$ 或 $2$，只有 $2$ 是素数只能是 $2$，域一定有特征，特征只能是 $2$。
\end{example}
